% lecture09.tex: Lecture 9: Rocky Planets I, Earth & Venus
% Planetary Systems, Kapteyn Institute

\documentclass[aspectratio=169, 11pt]{beamer}

% ── Theme and macros ────────────────────────────────────────
\usepackage{../common/beamerthemeIPS}
% macros.tex: Shared math macros for IPS lecture slides
% Mirrors definitions in book/_config.yml for consistency
% Uses \providecommand to avoid conflicts with unicode-math

% ── Derivatives ─────────────────────────────────────────────
\providecommand{\dv}[2]{\frac{\mathrm{d} #1}{\mathrm{d} #2}}
\providecommand{\dd}{}\renewcommand{\dd}{\mathrm{d}}
\providecommand{\pdv}[2]{\frac{\partial #1}{\partial #2}}

% ── Solar quantities ────────────────────────────────────────
\newcommand{\Msun}{M_{\odot}}
\newcommand{\Rsun}{R_{\odot}}
\newcommand{\Lsun}{L_{\odot}}

% ── Planetary quantities ────────────────────────────────────
\newcommand{\Mearth}{M_{\oplus}}
\newcommand{\Rearth}{R_{\oplus}}
\newcommand{\Mjup}{M_{\mathrm{J}}}
\newcommand{\Rjup}{R_{\mathrm{J}}}

% ── Constants ───────────────────────────────────────────────
\newcommand{\kB}{k_{\mathrm{B}}}

% ── Media links ─────────────────────────────────────────────
% Clickable button that opens the animated or video version of a
% figure, e.g. \mediabutton{https://...gif}{Play animation}
\providecommand{\mediabutton}[2]{\href{#1}{\beamergotobutton{#2}}}

\graphicspath{{figures/}{../common/}{../../book/09_earth_venus/figures/}}

% ── Metadata ────────────────────────────────────────────────
\title[Lecture 9: Earth \& Venus]{Lecture 9: Rocky Planets I\\Earth \& Venus}
\subtitle{Planetary Systems}
\author[L9]{Tim Lichtenberg}
\institute{Kapteyn Astronomical Institute\\University of Groningen}
\date{2026}
\titlebgcredit{Three Venus-bound missions over a Magellan radar mosaic. Credit: Widemann et al.\ (2023), SSRv 219, 56}

% ════════════════════════════════════════════════════════════
\begin{document}

% ── Title slide ─────────────────────────────────────────────
\begin{frame}[plain,noframenumbering]
  \titlepage
\end{frame}

% ── Learning objectives ─────────────────────────────────────
\begin{frame}{Learning Objectives}
  By the end of this lecture, you will be able to:
  \vskip8pt
  \begin{itemize}
    \item Use Earth's interior, surface, ocean and climate as the reference case
    \item Compare Venus to Earth in mass, radius, atmosphere and rotation
    \item Explain why a saturated water atmosphere has a maximum thermal flux
    \item Derive the Simpson-Nakajima runaway greenhouse limit
    \item Read the D/H ratio as evidence for past water on Venus
  \end{itemize}
\end{frame}

% ── Recap ───────────────────────────────────────────────────
\begin{frame}{Recap: Lectures 1--8}
  The two innermost large rocky planets share similar bulk parameters but followed radically divergent evolutionary paths.
  \vskip8pt
  \begin{itemize}
    \item Differentiation separates crust, mantle, and core ($C/MR^2 = 0.331$ for Earth)
    \item Volcanism, impact cratering, and escape shape surfaces and atmospheres
    \item Liquid outer core sustains a protective magnetic dynamo
  \end{itemize}
  \vskip8pt
  \textbf{Today:} What caused the evolutionary divergence of Earth and Venus?
\end{frame}


% ════════════════════════════════════════════════════════════
\sectionimage{pale_blue_dot}{Earth seen from $\sim 6$ billion km. Voyager 1 (1990). Credit: NASA/JPL-Caltech}
\section{Earth as Reference}
% ════════════════════════════════════════════════════════════

% ── Earth at a glance ───────────────────────────────────────
\begin{frame}{Earth at a Glance}
  \centering
  \small
  \begin{tabular}{l r l}
    \toprule
    \textbf{Parameter} & \textbf{Value} & \textbf{Note} \\
    \midrule
    Mass & $5.972 \times 10^{24}$ kg & defines $\Mearth$ \\
    Radius (mean) & $6371$ km & defines $\Rearth$ \\
    Mean density & $5515$ kg\,m$^{-3}$ & uncompressed $\sim 4030$ kg\,m$^{-3}$ \\
    Semimajor axis & $1.000$ AU & defines AU \\
    Bond albedo & $0.30$ & cloud + ice + land \\
    $T_{\mathrm{eq}}$ & $255$ K & equilibrium, no greenhouse \\
    $T_{\mathrm{surf}}$ & $288$ K & $+33$ K greenhouse warming \\
    Rotation period & $23.93$ h & spin axis tilt $23.4^\circ$ \\
    Magnetic field & dipolar, $\sim 25$--$65$ $\mu$T & active dynamo \\
    \bottomrule
  \end{tabular}
  \vskip6pt
  \sourcenote{NASA Earth Fact Sheet (NSSDC)}
\end{frame}

% ── Why is Earth habitable ──────────────────────────────────
\begin{frame}{What Makes Earth Habitable?}
  Four ingredients, all present and stable for $\sim 4$ Gyr:
  \vskip6pt
  \begin{enumerate}
    \item \textbf{Liquid water} on the surface, continuously
    \item \textbf{Plate tectonics}: long-term carbon cycle, climate buffering
    \item \textbf{Magnetic field}: shields atmosphere from solar wind erosion
    \item \textbf{Stable orbit}: low eccentricity, large stabilising Moon
  \end{enumerate}
  \vskip6pt
  \keyresult{Identify which planetary features are necessary, which are sufficient, and which are accidental.}
\end{frame}

% ── Earth in context: TRAPPIST-1 ────────────────────────────
\begin{frame}
  \centering
  \makebox[\linewidth][c]{\includegraphics[width=0.95\paperwidth, height=0.74\paperheight, keepaspectratio]{nasa_trappist1_solarsystem_comparison}}\\
  {\footnotesize TRAPPIST-1 vs.\ inner solar system, to scale. Seven Earth-sized worlds span the habitable zone around a late M dwarf, testing terrestrial evolutionary pathways. Credit: NASA/JPL-Caltech (PIA22093).}
\end{frame}

% ── Plate tectonics ─────────────────────────────────────────
\begin{frame}{Plate Tectonics: The Engine}
  Earth alone among known solar system bodies operates active plate tectonics today, driving continuous volatile recycling between surface and interior.
  \vskip8pt
  \begin{itemize}
    \item Lithosphere divided into $\sim 15$ major plates moving at cm\,yr$^{-1}$ speeds
    \item Subduction zones recycle crust and surface volatiles on $\sim 100$ Myr timescales
    \item Closes the long-term carbon cycle by returning buried carbonate to the mantle
  \end{itemize}
\end{frame}

% ── Geomagnetic field ───────────────────────────────────────
\begin{frame}{Earth's Magnetic Field}
  A vigorous geodynamo generates a protective dipolar magnetosphere that shields Earth's atmosphere from direct solar wind erosion.
  \vskip6pt
  \begin{columns}[c]
    \column{\recapfigcolnarrow}
    \ipsrecapfig{magnetosphere_anatomy_esa}
    \source{ESA, CC BY-SA 3.0 IGO}

    \column{\recaptextcolwide}
    \footnotesize
    \begin{itemize}
      \item Surface dipolar intensity of $25$--$65\ \mu\mathrm{T}$ driven by core convection
      \item Magnetopause dayside stand-off distance at $\sim 10\, \Rearth$
      \item Deflects solar wind and suppresses ion-pickup escape of volatiles
    \end{itemize}
  \end{columns}
\end{frame}

% ── Hydrosphere & cryosphere ────────────────────────────────
\begin{frame}{Hydrosphere and Cryosphere}
  Earth possesses a stable surface liquid ocean containing over a billion cubic kilometres of water, complemented by polar cryosphere reservoirs.
  \vskip8pt
  \begin{itemize}
    \item Oceans contain $1.34 \times 10^9\ \mathrm{km^3}$ of liquid water
    \item Total surface water equivalent to $\approx 2700$ m global equivalent layer (GEL)
    \item Deep mantle holds $\geq 1$--$2$ additional ocean masses
  \end{itemize}
\end{frame}

% ── Climate stability and solar evolution ───────────────────
\begin{frame}{Climate Stability and Solar Evolution}
  \begin{columns}[c]
    \column{\recapfigcol}
    \ipsrecapfig{zahnle2007_solar_evolution}
    \source{Zahnle et al.\ (2007)}

    \column{\recaptextcol}
    \textbf{Faint young Sun problem:}
    \begin{itemize}
      \item Solar luminosity increased by $\sim 30\%$ since $4.5$ Ga
      \item Liquid water persisted at $\geq 3.8$ Ga despite dimmer early Sun
      \item Carbonate-silicate negative feedback stabilized climate
    \end{itemize}
    \vskip4pt
    \keyresult{Habitable zone inner edge migrated outward over time, sweeping past Venus.}
  \end{columns}
\end{frame}


% ════════════════════════════════════════════════════════════
% ════════════════════════════════════════════════════════════


% ════════════════════════════════════════════════════════════
\sectionimage{venus_earth_mars_atmospheres}{Surface temperature and pressure for Venus, Earth, Mars. Course-original figure; data: NSSDC Planetary Fact Sheet}
\section{Venus: Overview and Surface}
% ════════════════════════════════════════════════════════════

% ── Earth-Venus systems context ─────────────────────────────
\begin{frame}{Earth and Venus as Coupled Systems}
  \begin{columns}[c]
    \column{\recapfigcol}
    \ipsrecapfig{gillmann2022_earth_venus_systems}
    \source{Gillmann et al.\ (2022)}

    \column{\recaptextcol}
    Both are coupled systems of:
    \begin{itemize}
      \item Interior + tectonics
      \item Magnetosphere or its absence
      \item Atmosphere + clouds
      \item Surface + (possible) hydrosphere
    \end{itemize}
    \vskip6pt
    The differences between them are emergent properties of which feedback loops close.
  \end{columns}
\end{frame}

% ── Venus comparison ────────────────────────────────────────
\begin{frame}{Venus and Earth: Twin and Not}
  \centering
  \small
  \begin{tabular}{l r r}
    \toprule
    \textbf{Parameter} & \textbf{Earth} & \textbf{Venus} \\
    \midrule
    Mass ($\Mearth$) & $1.000$ & $0.815$ \\
    Radius ($\Rearth$) & $1.000$ & $0.949$ \\
    Mean density (kg\,m$^{-3}$) & $5515$ & $5243$ \\
    Semimajor axis (AU) & $1.000$ & $0.723$ \\
    Stellar flux ($S_\oplus$) & $1.00$ & $1.91$ \\
    Bond albedo & $0.30$ & $0.77$ \\
    $T_{\mathrm{surf}}$ (K) & $288$ & $737$ \\
    Surface pressure (bar) & $1.0$ & $92$ \\
    Atmosphere & $\mathrm{N_2}$/$\mathrm{O_2}$ & $\mathrm{CO_2}$/$\mathrm{N_2}$ \\
    Rotation period (d) & $+1.0$ & $-243$ (retrograde) \\
    Magnetic field & $\sim 50$ $\mu$T & $< 10^{-5} \times$ Earth \\
    \bottomrule
  \end{tabular}
  \vskip6pt
  \sourcenote{NASA Venus Fact Sheet (NSSDC)}
\end{frame}

% ── Venus mission history ───────────────────────────────────
\begin{frame}{Venus: A Long Mission History}
  Hostile surface conditions of $92$ bar and $737$ K limited lander lifespans to two hours, focusing exploration on radar and orbital remote sensing.
  \vskip6pt
  \begin{itemize}
    \item \textbf{Venera} (1975--1982): first surface images and XRF basaltic composition
    \item \textbf{Magellan} (1990--1994): global radar mapping at $\sim 100$ m resolution
    \item \textbf{Venus Express \& Akatsuki} (2006 to present): atmospheric dynamics, escape rates (Venus Express) and super-rotation
  \end{itemize}
\end{frame}

% ── Venus surface terrains (hero) ───────────────────────────
\begin{frame}
  \centering
  \makebox[\linewidth][c]{\includegraphics[width=0.95\paperwidth, height=0.74\paperheight, keepaspectratio]{widemann2023_venus_geological_terrains}}\\
  {\footnotesize Magellan radar mosaic with terrain overlay. Volcanic plains cover $\sim$80\% of the surface alongside tessera highlands, showing no plate-tectonic boundaries. Widemann et al.\ (2023), \textit{SSRv} 219, 56.}
\end{frame}

% ── Crater age ──────────────────────────────────────────────
\begin{frame}{Crater Statistics: A Single Surface Age}
  A statistically random distribution of $\sim 1000$ craters yields a young global crater-retention age of 150--250 Myr across all terrains.
  \vskip6pt
  \begin{itemize}
    \item $\sim 1000$ catalogued impact craters ($\geq 2$ km) show no spatial clustering
    \item Uniform density implies a young surface age of $\sim 150$--$250$ Myr
    \item Debated between catastrophic global resurfacing and steady volcanic equilibrium
  \end{itemize}
\end{frame}

% ── Earth-Venus topography (hero) ───────────────────────────
\begin{frame}
  \centering
  \makebox[\linewidth][c]{\includegraphics[width=0.92\paperwidth, height=0.74\paperheight, keepaspectratio]{smrekar2018_earth_venus_topography}}\\
  {\footnotesize Topography and geoid comparison. Unimodal hypsometry on Venus contrasts with Earth's bimodal ocean-continent distribution, ruling out Earth-style plate tectonics. Smrekar et al.\ (2018).}
\end{frame}


% ════════════════════════════════════════════════════════════
\sectionimage{family_portrait}{Voyager 1 Solar System family portrait (1990); each of the six planets it detected appears in a small inset box. Credit: NASA/JPL-Caltech}
\section{Venus Interior and Tectonic Regime}
% ════════════════════════════════════════════════════════════

% ── Venus interior ──────────────────────────────────────────
\begin{frame}{Venus' Interior: Inferred, Not Measured}
  Without landed seismology or precession measurements, Venus' internal structure is modeled by Earth analogy rather than directly measured.
  \vskip6pt
  \begin{itemize}
    \item Moment of inertia factor $C/MR^2 \approx 0.33$ is assumed, not measured
    \item Inferred metallic core radius $R_c \approx 3000$--$3500$ km ($\sim 0.55\, R_V$)
    \item Mean density of $5243$ kg\,m$^{-3}$ requires a substantial iron core
  \end{itemize}
\end{frame}

% ── Tectonic evolution schematic ───────────────────────────
\begin{frame}{Venus Tectonic Evolution: A Model Trajectory}
  \begin{columns}[c]
    \column{\recapfigcolnarrow}
    \ipsrecapfig{smrekar2018_venus_tectonic_evolution}
    \source{Smrekar et al.\ (2018)}

    \column{\recaptextcolwide}
    Surface temperature evolution from one numerical model: successive intervals of
    \begin{itemize}
      \item \textbf{Stagnant lid}
      \item \textbf{Mobile lid}
      \item \textbf{Stagnant lid}
      \item \textbf{Episodic lid}
    \end{itemize}
    \vskip4pt
    Conceptual alternatives debated for Venus today: \textbf{steady stagnant} (localised mantle plumes) vs.\ \textbf{episodic catastrophic overturn}. Either way, no Earth-style plate tectonics.
  \end{columns}
\end{frame}

% ── No magnetic field ───────────────────────────────────────
\begin{frame}{The Missing Magnetic Field}
  Venus lacks an intrinsic dipole magnetic field, leaving its upper atmosphere exposed to direct solar wind interaction and erosion.
  \vskip6pt
  \begin{itemize}
    \item Internal magnetic field measured at $< 10^{-5}$ times Earth's field
    \item Stagnant-lid insulation suppresses core heat loss, shutting down thermal convection
    \item Unshielded atmosphere enables solar-wind induced escape and ion pickup
  \end{itemize}
\end{frame}

% ── Stagnant lid vs episodic ────────────────────────────────
\begin{frame}{Why Stagnant Lid?}
  A hot, dry surface produces a strong lithosphere that resists subduction, locking Venus into an immobile stagnant-lid convective regime.
  \vskip6pt
  \begin{itemize}
    \item Plate subduction requires a hydrated, mechanically weakened lithosphere
    \item High surface temperatures ($737$ K) promote ductile flow rather than brittle faulting
    \item No subduction return leg breaks the long-term carbon cycle permanently
  \end{itemize}
  \vskip4pt
  \sourcenote{Smrekar et al.\ (2018); Gillmann et al.\ (2022)}
\end{frame}

% ── Recent volcanism ────────────────────────────────────────
\begin{frame}{Is Venus Volcanically Active Today?}
  Radar and orbital observations demonstrate that Venus remains volcanically active today, driven by internal heat loss through mantle plumes.
  \vskip6pt
  \begin{itemize}
    \item Magellan radar reanalysis reveals a $\sim 2\ \mathrm{km^2}$ vent change on Maat Mons (1991--1992)
    \item Venus Express thermal emissivity anomalies indicate lava flows younger than $250$ kyr
    \item Decadal fluctuations in atmospheric $\mathrm{SO_2}$ suggest episodic volcanic outgassing
  \end{itemize}
\end{frame}


% ════════════════════════════════════════════════════════════
\sectionimage{venus_earth_mars_atmospheres}{Surface temperature and pressure for Venus, Earth, Mars. Course-original figure; data: NSSDC Planetary Fact Sheet}
\section{Venus Atmosphere and Climate}
% ════════════════════════════════════════════════════════════

% ── Venus atmosphere structure ──────────────────────────────
\begin{frame}{Venus Atmosphere: Structure}
  Venus hosts a massive 92-bar carbon dioxide atmosphere enveloped in global sulfuric acid clouds without a tropopause cold trap.
  \vskip6pt
  \begin{itemize}
    \item Composition: $96.5\%\ \mathrm{CO_2}$, $3.5\%\ \mathrm{N_2}$ at $737$ K and $92$ bar surface pressure
    \item Global cloud deck at $48$--$70$ km composed of $\sim 75\%$ sulfuric acid droplets
    \item Lower atmosphere follows a dry adiabat with lapse rate $\Gamma \approx 8\ \mathrm{K\,km^{-1}}$; scale height $H \approx 15.9$ km
    \item $^{40}\mathrm{Ar}/^{36}\mathrm{Ar} \approx 1.1$ (Earth $\sim 300$): Venus has degassed only a fraction of its interior
  \end{itemize}
\end{frame}

% ── Super-rotation ──────────────────────────────────────────
\begin{frame}{Atmospheric Super-Rotation}
  Venus exhibits atmospheric super-rotation, where high-altitude cloud winds encircle the planet sixty times faster than the underlying solid surface.
  \vskip6pt
  \begin{itemize}
    \item Cloud-top winds of $\sim 100\ \mathrm{m\,s^{-1}}$ circle the planet in $\sim 4$ days
    \item Solid-body retrograde rotation takes $243$ days ($60\times$ angular velocity ratio)
    \item Maintained by wave transport and atmospheric thermal tides
  \end{itemize}
\end{frame}

% ── Greenhouse warming ──────────────────────────────────────
\begin{frame}{Venus' Greenhouse: $T_{\mathrm{surf}}$ from $T_{\mathrm{eq}}$}
  Despite absorbing less solar energy than Earth due to high albedo, Venus experiences extreme greenhouse warming from its massive carbon dioxide atmosphere.
  \vskip4pt
  $$T_{\mathrm{eq}} = \left[\frac{S(1-A)}{4\sigma}\right]^{1/4} \approx 227\ \mathrm{K},\quad T_{\mathrm{surf}} = 737\ \mathrm{K}$$
  \begin{itemize}
    \item Bond albedo $A = 0.77$ reflects most sunlight, yielding $T_{\mathrm{eq}} = 227\ \mathrm{K}$
    \item One-layer greenhouse $T_s = 2^{1/4}\,T_{\mathrm{eq}} \approx 270\ \mathrm{K}$ falls far short
    \item Greenhouse warming $\Delta T = +510\ \mathrm{K}$ dwarfs Earth's $+33\ \mathrm{K}$
    \item Driven by $92$ bar of $\mathrm{CO_2}$, with $\mathrm{SO_2}$, $\mathrm{H_2O}$ and the clouds closing the spectral windows
  \end{itemize}
\end{frame}

% ── Cloud emission spectrum ─────────────────────────────────
\begin{frame}{Where the Night-Side Clouds Form}
  \begin{columns}[c]
    \column{\recapfigcolwide}
    \ipsrecapfig{turbet2021_water_clouds_emission}
    \source{Turbet et al.\ (2021), Fig.~2a-d}

    \column{\recaptextcolnarrow}
    \footnotesize
    \begin{itemize}
      \item Hot and steamy Earth (a, c) and Venus (b, d) at the same flux, 340.5~W/m$^2$
      \item Clouds gather on the night-side in both, although Venus turns 243 times slower: the pattern does not need slow rotation
      \item Night-side clouds trap thermal radiation, so the emission maps are anticorrelated with the cloud maps
    \end{itemize}
  \end{columns}
\end{frame}

% ── Runaway threshold concept (hero) ────────────────────────
\begin{frame}
  \centering
  \makebox[\linewidth][c]{\includegraphics[width=0.88\paperwidth, height=0.74\paperheight, keepaspectratio]{zahnle2007_runaway_threshold}}\\
  {\footnotesize Thermal radiation limit and runaway greenhouse threshold. A saturated atmosphere has a maximum OLR of $\sim 280$--$300\ \mathrm{W\,m^{-2}}$, beyond which surface water vaporises completely. Zahnle et al.\ (2007).}
\end{frame}


% ── Break ───────────────────────────────────────────────────
\breakslide


% ════════════════════════════════════════════════════════════
\sectionimage{family_portrait}{Voyager 1 Solar System family portrait (1990); each of the six planets it detected appears in a small inset box. Credit: NASA/JPL-Caltech}
\section{Blackboard Derivation: Simpson-Nakajima Limit}
% ════════════════════════════════════════════════════════════

% ── Setup ───────────────────────────────────────────────────
\begin{frame}{Setup: Saturated Atmosphere, Greybody}
  \textbf{Goal:} show that a water-saturated atmosphere has a maximum OLR independent of surface temperature.
  \vskip6pt
  \textbf{Assumptions:}
  \begin{itemize}
    \item Atmosphere saturated in water vapour everywhere
    \item Single grey IR opacity $\kappa$, units m$^2$\,kg$^{-1}$
    \item Hydrostatic equilibrium, surface gravity $g$
    \item Photosphere: where IR optical depth $\tau \sim 1$
  \end{itemize}
  \vskip8pt
  Clausius-Clapeyron sets the water vapour partial pressure:
  $$\dv{p_{\mathrm{sat}}}{T} = \frac{L\, p_{\mathrm{sat}}}{R_v\, T^2}$$
  with $L \approx 2.5 \times 10^6$ J\,kg$^{-1}$, $R_v = 461$ J\,kg$^{-1}$\,K$^{-1}$.
\end{frame}

% ── Step 3: OLR asymptote ───────────────────────────────────
\begin{frame}{Result: The OLR Asymptote}
  Optical depth unity fixes the photosphere pressure, saturation fixes its temperature (board):
  $$p_{\mathrm{phot}} = \frac{g}{\kappa} = p_{\mathrm{sat}}(T_{\mathrm{phot}}) \quad\Rightarrow\quad T_{\mathrm{phot}} \approx 260\ \mathrm{K}$$
  $$F_{\mathrm{OLR}}^{\max} \approx \sigma\, T_{\mathrm{phot}}^4 \approx 5.67\times 10^{-8}\times (260)^4 \approx 260\, \mathrm{W\,m^{-2}}$$
  \vskip4pt
  \keyresult{Refined line-by-line calculations: $F_{\mathrm{OLR}}^{\max} \approx 280$--$310$ W\,m$^{-2}$ (Kasting 1988; Goldblatt et al.\ 2013).}
  \vskip4pt
  As $T_s$ rises, the photosphere migrates upward but its temperature stays clamped: thermal flux saturates.
\end{frame}

% ── Absorbed flux against the limit ─────────────────────────
\begin{frame}{Against the Limit: Earth Below, Venus Above}
  $$F_{\mathrm{abs}} = \frac{S_\odot}{4 a^2}\,(1 - A)$$
  \begin{itemize}
    \item Earth ($a = 1$ AU, $A = 0.30$): $F_{\mathrm{abs}} \approx 240$ W\,m$^{-2}$, below $F_{\mathrm{OLR}}^{\max}$
    \item Early Venus ($a = 0.723$ AU, Earth-like $A$): $F_{\mathrm{abs}} \approx 460$ W\,m$^{-2}$, far above it
    \item Present Venus ($A = 0.77$) absorbs only $\sim 150$ W\,m$^{-2}$: the albedo is the product of the runaway, not its cause
  \end{itemize}
  \vskip4pt
  \keyresult{Setting $F_{\mathrm{abs}} = F_{\mathrm{OLR}}^{\max}$ gives the inner edge of the habitable zone, $a_{\mathrm{IHZ}} \approx 0.95$--$0.99$ AU.}
\end{frame}

% ── Kopparapu inner HZ (hero) ───────────────────────────────
\begin{frame}
  \centering
  \makebox[\linewidth][c]{\includegraphics[width=0.90\paperwidth, height=0.74\paperheight, keepaspectratio]{kopparapu2013_inner_HZ}}\\
  {\footnotesize Inner habitable zone calculation. The runaway greenhouse limits the inner edge to $S_{\mathrm{eff}} = 1.06$ ($0.97$ AU around the present Sun), placing Venus ($1.91\,S_\oplus$) well inside. Kopparapu et al.\ (2013).}
\end{frame}


% ════════════════════════════════════════════════════════════
\sectionimage{terrestrial_planets}{Mercury, Venus, Earth, Mars at approximate relative scale. Credit: NASA/JPL}
\section{Venus Water Loss and the D/H Ratio}
% ════════════════════════════════════════════════════════════

% ── When did water leave ────────────────────────────────────
\begin{frame}{When Did Venus Lose Its Water?}
  Two hypotheses: early loss from a magma ocean, or late desiccation of a temperate ocean.
  \vskip6pt
  \begin{itemize}
    \item \textbf{Early loss (Hamano Type II):} the steam atmosphere never condenses; water escapes during the magma-ocean phase
    \item \textbf{Late loss (Way et al.\ 2016):} a temperate surface persists for Gyr until solar brightening triggers the runaway
    \item The timing needs noble-gas isotopes and surface mineralogy from the next missions
  \end{itemize}
\end{frame}

% ── Way 2016 paleo-Venus ────────────────────────────────────
\begin{frame}{Was Venus Ever Habitable? The Late-Loss Scenario}
  \begin{columns}[c]
    \column{\recapfigcol}
    \ipsrecapfig{way2016_paleo_venus_temperature}
    \source{Way et al.\ (2016)}

    \column{\recaptextcol}
    \begin{itemize}
      \item GCM snapshot at 2.9 Gya: surface $T$ ranges from $-22^{\circ}$C polar to $+46^{\circ}$C equatorial
      \item Cloud feedback from slow rotation reflects flux back to space
      \item Across the Way et al.\ (2016) ensemble, global mean $T < 290$ K can persist for $\sim 1$ Gyr
      \item A habitable Venus is dynamically possible, not required
    \end{itemize}
  \end{columns}
\end{frame}

% ── Photolysis + H escape ───────────────────────────────────
\begin{frame}{The Loss Mechanism: Photolysis + Escape}
  Massive water loss proceeds sequentially via stratospheric photolysis, rapid hydrogen escape to space, and residual oxygen oxidation or ion pickup.
  \vskip6pt
  \begin{itemize}
    \item UV photolysis ($\mathrm{H_2O} + h\nu \rightarrow \mathrm{H} + \mathrm{OH}$) dissociates water in the upper atmosphere
    \item Light hydrogen escapes efficiently through thermal Jeans and hydrodynamic escape
    \item Oxygen is absorbed by crustal oxidation or stripped by solar-wind ion pickup
  \end{itemize}
\end{frame}

% ── D/H ratio ───────────────────────────────────────────────
\begin{frame}{The D/H Ratio: Evidence for Lost Water}
  Venus' atmospheric D/H is $\sim 150\times$ Earth's: the record of massive past hydrogen escape.
  \vskip6pt
  \begin{columns}[c]
    \column{0.52\textwidth}
    \centering
    \small
    \begin{tabular}{l r}
      \toprule
      \textbf{Body} & \textbf{D/H ($\times 10^{-4}$)} \\
      \midrule
      Protosolar (H$_2$) & $\sim 0.20$ \\
      Earth (oceans) & $1.56$ \\
      Comets (Halley, 67P) & $1.5$--$5$ \\
      \textbf{Venus} & $\sim \mathbf{240}$ \\
      Mars & $\sim 8$ \\
      \bottomrule
    \end{tabular}
    \vskip4pt
    \sourcenote{Donahue et al.\ (1982); de Bergh et al.\ (1991)}

    \column{0.46\textwidth}
    \begin{itemize}
      \item Light $^1\mathrm{H}$ escapes more easily than $^2\mathrm{H}$
      \item Rayleigh distillation: $R/R_0 = f^{\alpha-1}$
      \item Diffusion-limited ($\alpha = 0.5$): $f \approx 150^{-2}$, initial water $\sim 450$ m GEL; the Jeans limit ($\alpha = 0.71$) gives an unphysical $> 100$ oceans, so 450 m is a lower bound
    \end{itemize}
  \end{columns}
\end{frame}

% ── Initial water inventory ─────────────────────────────────
\begin{frame}{Bounding the Initial Water Inventory}
  Rayleigh distillation models translate the measured deuterium enrichment into a firm lower bound on Venus' initial water inventory.
  \vskip6pt
  \begin{itemize}
    \item Rayleigh distillation relation: $R/R_0 = f^{\alpha - 1}$ with fractionation factor $\alpha$
    \item Diffusion limit ($\alpha \approx 0.5$) yields remaining water fraction $f \approx 150^{-2} \approx 4.4 \times 10^{-5}$
    \item Current $2$ cm GEL column implies initial water equivalent layer $d_{\mathrm{init}} \geq 450$ m GEL
  \end{itemize}
\end{frame}

% ── Turbet hysteresis ───────────────────────────────────────
\begin{frame}{Bistability and Hysteresis Loop}
  \begin{columns}[c]
    \column{\recapfigcol}
    \ipsrecapfig{turbet2021_hysteresis}
    \source{Turbet et al.\ (2021), Fig.~4}

    \column{\recaptextcol}
    For some stellar fluxes, two stable states coexist:
    \begin{itemize}
      \item Wet: ocean + thin atmosphere
      \item Dry: hot, thick $\mathrm{CO_2}$ atmosphere
    \end{itemize}
    \vskip6pt
    Crossing the threshold is one-way. The dry state has no thermodynamic path back.
  \end{columns}
\end{frame}

% ── Constantinou Venus pathways ─────────────────────────────
\begin{frame}{Constraining the Trajectory}
  \begin{columns}[c]
    \column{\recapfigcolnarrow}
    \ipsrecapfig{constantinou2024_venus_pathways}
    \source{Constantinou et al.\ (2024)}

    \column{\recaptextcolwide}
    Two end-member pathways for Venus' water history:
    \begin{itemize}
      \item \textbf{Top branch (dry):} no surface ocean; H lost early; mantle dry today
      \item \textbf{Bottom branch (wet):} ocean condenses, late loss; mantle still wet today
      \item Constantinou et al.\ (2024) argue current photochemistry favours the dry/early-loss branch (volcanic $\mathrm{H_2O} \lesssim 6\%$)
      \item DAVINCI noble-gas measurements will discriminate
    \end{itemize}
  \end{columns}
\end{frame}


% ════════════════════════════════════════════════════════════
\sectionimage{terrestrial_planets}{Mercury, Venus, Earth, Mars at approximate relative scale. Credit: NASA/JPL}
\section{Why Did Earth and Venus Diverge?}
% ════════════════════════════════════════════════════════════

% ── Carbonate-silicate failure ──────────────────────────────
\begin{frame}{Carbonate-Silicate Cycle: Failure Mode}
  Without surface liquid water and subduction, Venus lacks weathering sinks and mantle return legs, causing volcanic carbon dioxide to build indefinitely.
  \vskip6pt
  \begin{itemize}
    \item Silicate weathering sink collapses entirely in the absence of liquid surface water
    \item Stagnant-lid regime prevents subduction return of carbonates to the mantle
    \item Volcanic outgassing irreversibly drives atmospheric $\mathrm{CO_2}$ to the observed $92$ bar
  \end{itemize}
\end{frame}

% ── Lammer carbonate-silicate ───────────────────────────────
\begin{frame}{The Earth Carbonate-Silicate Schematic}
  \begin{columns}[c]
    \column{\recapfigcol}
    \ipsrecapfig{lammer2018_carbonate_silicate}
    \source{Lammer et al.\ (2018)}

    \column{\recaptextcol}
    Earth's closed carbon loop:
    \begin{itemize}
      \item Volcanic outgassing (mantle to atmosphere)
      \item Silicate weathering on continents (atmosphere to ocean, $T$-dependent sink)
      \item Carbonate sedimentation + subduction (ocean to mantle return leg)
    \end{itemize}
    \vskip4pt
    Venus today has no liquid-water sink and no subduction return: both halves of the loop are broken.
  \end{columns}
\end{frame}

% ── Lammer accretion ────────────────────────────────────────
\begin{frame}{Initial Volatile Inventories}
  \begin{columns}[c]
    \column{\recapfigcol}
    \ipsrecapfig{lammer2018_venus_earth_mars_accretion}
    \source{Lammer et al.\ (2018)}

    \column{\recaptextcol}
    \begin{itemize}
      \item Comparative accretion of Venus, Earth, Mars
      \item Initial water + carbon inventories overlap substantially
      \item Divergence is not from initial conditions alone
      \item Early evolution decides which branch each planet takes
    \end{itemize}
  \end{columns}
\end{frame}

% ── Divergence and habitability implications ────────────────
\begin{frame}{Why the Twins Diverged: Habitability Implications}
  Coupled factors drove Earth and Venus apart:
  \vskip6pt
  \begin{enumerate}
    \item \textbf{Solar flux:} $1.91\times$ Earth's insolation, above the runaway threshold
    \item \textbf{Geophysics:} no intrinsic dynamo, slow retrograde rotation
    \item \textbf{Tectonics:} a stagnant lid removes the subduction return leg of the carbon cycle
  \end{enumerate}
  \vskip6pt
  \keyresult{Venus defines the inner edge of the habitable zone: coupled interior-atmosphere feedbacks decide whether a terrestrial planet keeps its water.}
\end{frame}


% ════════════════════════════════════════════════════════════
\sectionimage{widemann2023_three_missions_render}{VERITAS, DAVINCI and EnVision over an Akatsuki UV view of Venus. Credit: Widemann et al.\ (2023)}
\section{Recent Advances and Future Missions}
% ════════════════════════════════════════════════════════════

% ── Three missions hero ─────────────────────────────────────
\begin{frame}
  \centering
  \makebox[\linewidth][c]{\includegraphics[width=0.95\paperwidth, height=0.74\paperheight, keepaspectratio]{widemann2023_three_missions_render}}\\
  {\footnotesize The 2030s Venus fleet: DAVINCI (NASA), VERITAS (NASA), and EnVision (ESA). DAVINCI's descent probe mass spectrometer will measure atmospheric noble gases to distinguish early against late water loss. Widemann et al.\ (2023).}
\end{frame}

% ── Summary ─────────────────────────────────────────────────
\begin{frame}{Summary}
  \begin{enumerate}
    \item \textbf{Bulk twin:} same size and density, but $\sim 500$ K and $92$ bar apart
    \item \textbf{Runaway limit:} the Simpson-Nakajima threshold ($280$--$310\ \mathrm{W\,m^{-2}}$) sets the inner habitable zone
    \item \textbf{D/H evidence:} $\sim 150\times$ Earth's ratio records massive water loss
    \item \textbf{Broken cycle:} no surface water stops the weathering sink; a stagnant lid removes the subduction return leg
  \end{enumerate}
  \vskip8pt
  \textbf{Next: Lecture 10} Mercury and Mars.
\end{frame}

% ── Final slide ─────────────────────────────────────────────
\begin{frame}[plain,noframenumbering]
  \begin{tikzpicture}[remember picture, overlay]
    \ipscontain{title_background}
    \fill[ipsPrimary, opacity=0.70] (current page.north west) rectangle (current page.south east);
  \end{tikzpicture}
  \vfill
  \begin{center}
    {\color{white}\Large\bfseries Questions?}
    \vskip20pt
    {\color{ipsLightFill}\normalsize Next: \textbf{Lecture~10. Rocky Planets II: Mercury \& Mars}}
    \vskip16pt
    {\color{ipsLightFill!70}\small Lecture notes: \texttt{formingworlds.github.io/IntroductionToPlanetaryScience}}
  \end{center}
  \vfill
\end{frame}

\end{document}
